\documentclass[11pt]{article}
\input{../../shared/project_style.tex}
\rhead{Basics 26 Textbook Draft}

\begin{document}
\DocTitle{Basics 26: Ideal Ohm's Law and Flux Freezing}{IV - Magnetohydrodynamics}{The moving-fluid electric field, induction, Alfvén's theorem, topology, and nonideal breakdown}

\begin{overviewbox}
\textbf{Textbook draft, version 1.0.} The relation $\mathbf E+\mathbf u\times\mathbf B=0$ is the kinematic core of ideal MHD. This chapter derives its interpretation, proves magnetic-flux conservation for a moving surface, and explains precisely how resistive and kinetic effects break ideal topology. This chapter is designed for a reader with no prior plasma-physics training. It distinguishes exact identities from physical approximations and connects the central equations to later simulation modules.
\end{overviewbox}

\ProjectTOC

\section{Learning objectives}
After completing this note, the reader should be able to:
\begin{itemize}[leftmargin=1.6em]
\item Interpret ideal Ohm's law in the fluid rest frame.
\item Derive the ideal induction equation.
\item Derive the field-to-density stretching relation.
\item Prove flux freezing for a material surface.
\item Distinguish flux conservation from constant field strength.
\item Identify physical and numerical mechanisms that break flux freezing.
\end{itemize}

\section{Prerequisites and reading strategy}
\begin{itemize}[leftmargin=1.6em]
\item Basics 3, 12, 21, and 24.
\end{itemize}
On a first reading, emphasize physical meaning and dimensional checks. On a second reading, reproduce the derivations. Attempt the computational laboratory only after the limiting cases are understood.

\section{Electric field in the moving fluid frame}

At nonrelativistic speeds, the leading transformation to a frame moving with $\mathbf u$ is
\begin{equation}
\mathbf E'=\mathbf E+\mathbf u\times\mathbf B.
\end{equation}
Ideal Ohm's law,
\begin{equation}
\boxed{\mathbf E+\mathbf u\times\mathbf B=0},
\end{equation}
states that the electric field vanishes in the local bulk-fluid frame. It does not state that the laboratory electric field vanishes.


\section{From generalized to ideal Ohm's law}

Electron momentum balance gives schematically
\begin{equation}
\mathbf E+\mathbf u\times\mathbf B
=\eta\mathbf J+\frac{\mathbf J\times\mathbf B}{en_e}
-\frac{\nabla\cdot\mathsf P_e}{en_e}
-\frac{m_e}{e}\frac{D_e\mathbf u_e}{Dt}.
\end{equation}
Ideal MHD orders every term on the right small. The condition is therefore scale dependent. A plasma can be ideal globally while containing localized nonideal layers.


\section{Induction and magnetic-field stretching}

Faraday's law gives
\begin{equation}
\partial_t\mathbf B=\nabla\times(\mathbf u\times\mathbf B).
\end{equation}
Use
\begin{equation}
\nabla\times(\mathbf u\times\mathbf B)=
\mathbf u(\nabla\cdot\mathbf B)-\mathbf B(\nabla\cdot\mathbf u)
+(\mathbf B\cdot\nabla)\mathbf u-(\mathbf u\cdot\nabla)\mathbf B.
\end{equation}
With $\nabla\cdot\mathbf B=0$,
\begin{equation}
\frac{D\mathbf B}{Dt}=(\mathbf B\cdot\nabla)\mathbf u-\mathbf B\nabla\cdot\mathbf u.
\end{equation}
Combining with $D\rho/Dt=-\rho\nabla\cdot\mathbf u$ gives
\begin{equation}
\boxed{\frac{D}{Dt}\left(\frac{\mathbf B}{\rho}\right)
=\left(\frac{\mathbf B}{\rho}\cdot\nabla\right)\mathbf u}.
\end{equation}
The field-to-mass ratio is stretched by velocity gradients.


\section{Flux-freezing theorem}

Let $S(t)$ move with the fluid. The transport theorem for magnetic flux gives
\begin{equation}
\frac{d}{dt}\int_{S(t)}\mathbf B\cdot d\mathbf A
=\int_{S(t)}
\left[\partial_t\mathbf B-\nabla\times(\mathbf u\times\mathbf B)\right]\cdot d\mathbf A.
\end{equation}
The bracket vanishes by the ideal induction equation:
\begin{equation}
\boxed{\frac{d\Phi_B}{dt}=0}.
\end{equation}
Equivalently, Faraday's law plus the motional electromotive force around a material loop gives zero net change.


\section{Physical interpretation and cautions}

If two ideal-fluid elements lie on the same magnetic field line initially, the mapping keeps them connected under smooth evolution. Field lines may be compressed, stretched, and folded, so $B$ can change greatly even though material-surface flux is fixed. The statement is topological and Lagrangian; it is not the assertion that drawn field lines are physical strings embedded in matter.


\section{Breaking ideal flux freezing}

With resistivity,
\begin{equation}
\partial_t\mathbf B=\nabla\times(\mathbf u\times\mathbf B)+\eta_m\nabla^2\mathbf B
\end{equation}
for uniform magnetic diffusivity $\eta_m=\eta/\mu_0$. The magnetic Reynolds number
\begin{equation}
R_m=\frac{UL}{\eta_m}
\end{equation}
compares advection with diffusion. Hall, pressure-tensor, inertia, finite-Larmor-radius, and turbulent effects can also break the bulk-fluid frozen-in relation, often with different species carrying distinct field-line velocities.


\section{Numerical flux freezing}

A discrete method may change topology because of truncation error even when the PDE is ideal. Excess numerical resistivity can damp Alfvén waves and alter eigenmode continua. Diagnostics include magnetic-energy decay, flux through advected loops, and convergence of reconnection rates toward zero in an ideal benchmark.


\section{Computational laboratory}
Advect a two-dimensional divergence-free magnetic field with a prescribed incompressible velocity. Track flux through a moving material loop and compare centered, upwind, and constrained-transport discretizations. Add controlled resistivity and recover the expected magnetic-diffusion scaling.

\section{Summary}
\begin{itemize}[leftmargin=1.6em]
\item Ideal Ohm's law means the electric field vanishes in the bulk-fluid frame.
\item It produces induction by advection, compression, and stretching of magnetic field.
\item Magnetic flux through a material surface is conserved under smooth ideal evolution.
\item Resistive, Hall, pressure, inertia, and kinetic effects can break the bulk frozen-in condition.
\item Numerical diffusion can mimic nonideal physics and must be quantified.
\end{itemize}

\SymbolsAcronymsSection
\begin{symtable}
$\mathbf E'$ & moving-frame electric field & $\mathbf E+\mathbf u\times\mathbf B$. \\
$\Phi_B$ & magnetic flux & Surface integral of magnetic field. \\
$\eta_m$ & magnetic diffusivity & $\eta/\mu_0$. \\
$R_m$ & magnetic Reynolds number & Advection-to-diffusion ratio. \\
$S(t)$ & material surface & Surface moving with the fluid. \\
$D/Dt$ & material derivative & $\partial_t+\mathbf u\cdot\nabla$. \\
\end{symtable}

\section{Exercises}
\begin{enumerate}[leftmargin=1.8em]
\item Derive the expanded material form of the induction equation.
\item Combine mass and induction equations to obtain the evolution of $\mathbf B/\rho$.
\item Prove flux freezing using the moving-surface transport theorem.
\item Explain why magnetic-field strength can change under frozen flux.
\item Estimate the magnetic Reynolds number for a representative fusion-plasma scale.
\end{enumerate}

\section{References}
\begin{thebibliography}{99}
\bibitem{ref1} D. J. Griffiths, \emph{Introduction to Electrodynamics}, 4th ed., Cambridge University Press (2017).
\bibitem{ref2} J. D. Jackson, \emph{Classical Electrodynamics}, 3rd ed., Wiley (1998).
\bibitem{ref3} F. F. Chen, \emph{Introduction to Plasma Physics and Controlled Fusion}, 3rd ed., Springer (2016).
\bibitem{ref4} R. J. Goldston and P. H. Rutherford, \emph{Introduction to Plasma Physics}, CRC Press (1995).
\bibitem{ref5} P. M. Bellan, \emph{Fundamentals of Plasma Physics}, Cambridge University Press (2006).
\bibitem{ref6} J. P. Freidberg, \emph{Ideal MHD}, Cambridge University Press (2014).
\bibitem{ref7} J. P. Goedbloed, R. Keppens, and S. Poedts, \emph{Advanced Magnetohydrodynamics}, Cambridge University Press (2010).
\bibitem{ref8} H. Goedbloed and S. Poedts, \emph{Principles of Magnetohydrodynamics}, Cambridge University Press (2004).
\end{thebibliography}

\end{document}
